% =====================================================================
% Deep literature review: knowledge-graph-enhanced and hybrid scholarly
% search.  Assumes standard packages (hyperref, natbib or biblatex)
% and a bibliography file containing the keys used below.
% =====================================================================

\subsection{Recent research advances}
\label{sec:recent-advances}

This thesis takes a different stance toward each of the four threads
examined below.  Dense retrieval (Section~\ref{sec:lr-dense}) is retained
as one signal in a hybrid pipeline, since the strongest current models
still miss roughly a quarter of relevant papers on realistic scientific
queries.  Graph-augmented methods (Section~\ref{sec:lr-graph}) supply the
conceptual motivation for ingesting OpenAlex into a typed knowledge
graph but offer no benchmarked baseline to copy.  Structured-query
translation (Section~\ref{sec:lr-nl2q}) is the design this thesis
pursues, instantiated through an intermediate JSON representation that
decouples LLM-based intent extraction from deterministic query
compilation.  Reformulation (Section~\ref{sec:lr-rag}) is set aside, both
because it degrades strong retrievers and because it cannot express the
structural queries a typed graph makes available.

\subsubsection{Scholarly embeddings and dense retrieval for scientific text}
\label{sec:lr-dense}

Dense retrieval is retained as one signal in this thesis's hybrid
pipeline rather than its primary mechanism, because the strongest current
models still miss roughly a quarter of relevant papers on realistic
scientific queries.  Three lines of progress have improved dense
retrieval without closing this gap.  Wider training distributions, from
DPR \cite{karpukhin2020dpr} through unsupervised contrastive
pre-training \cite{izacard2022contriever} and weakly-supervised pairs at
$10^8$--$10^9$ scale \cite{wang2022e5,li2023gte,xiao2023cpack,chen2024bgem3},
have lifted average nDCG@10 on BEIR \cite{thakur2021beir} but left a
structural weakness on entity-heavy and lexical queries.  Late
interaction \cite{khattab2020colbert,santhanam2022colbertv2} and learned
sparse retrieval \cite{formal2021splade,formal2022splade} preserved
fine-grained query-document interaction at modest cost, with ColBERTv2
reaching a BEIR nDCG@10 average of 49.9 across thirteen datasets while
shrinking the index by a factor of 6--10 \cite{santhanam2022colbertv2}.
Scholarly-domain encoders pursued the orthogonal axis of citation
supervision \cite{beltagy2019scibert,cohan2020specter,ostendorff2022scincl},
but Singh et al.'s 24-task SciRepEval benchmark
\cite{singh2023scirepeval} shows that proximity-trained representations
do not transfer to classification or ad-hoc search.  Their adapter-based
SPECTER2 reaches MAP 38.4 and Recall@5 33.0 on MDCR against BM25's 33.7
and 28.5, suggesting single-vector citation-trained encoders are not
sufficient for heterogeneous scholarly tasks.

The strongest empirical evidence that dense retrieval remains unsolved
on scientific text is LitSearch \cite{ajith2024litsearch}.  Ajith et al.\
assembled 597 realistic literature-search questions paired against a
corpus of 64,183 papers, half rewritten from inline citations and half
written by the papers' own authors.  None of the systems perform well.
The best open model tested, GritLM-7B, reaches Recall@5 of 74.8 on
specific questions but falls to 67.7 on the inline-citation subset, which
deliberately filters for low title-abstract lexical overlap, and to
Recall@20 of 70.8 on broad questions.  Commercial tools score 17.5
Recall@5 on author-written specific questions, roughly 65 points below
GritLM-7B.  The cross-benchmark contrast makes the diagnosis concrete:
GritLM-7B scores 70.3 nDCG@10 on Natural Questions but only 44.1 on broad
LitSearch, a 26-point drop that cannot be explained by random noise.
LitSearch also surfaces three failure modes that single-vector dense
retrieval cannot easily address.  Inline-citation queries are
systematically harder than author-written ones because authors reuse
their own title and abstract vocabulary, and retrievers exploit this
lexical shortcut.  Broad questions remain harder than specific ones.
Feeding full text rather than title and abstract to GritLM hurts
retrieval on author-written specific queries (Recall@5 drops from 82.5 to
73.0), which Ajith et al.\ attribute to a distributional mismatch between
training (MS~MARCO passages around 60 tokens) and full scientific texts.
Related work reinforces the picture.  DORIS-MAE \cite{wang2023dorismae}
reports that general-purpose models modestly outperform most
scholarly-specific ones on multi-aspect CS queries, and CSFCube
\cite{mysore2021csfcube} shows that SPECTER's performance varies
substantially across facets.  Notably, LitSearch does not benchmark
SPECTER2 or SciNCL in its main tables.  Combined with the DORIS-MAE
findings, this strongly suggests that citation-triplet-trained encoders
have not been shown to outperform strong general-purpose retrievers on
realistic natural-language literature search.  They remain the models of
choice for paper-paper similarity, but that is a different task.

These structural failures, on broad questions, multi-hop intent, and
typed-relation queries that single-vector representations cannot express,
are precisely what motivates adding graph traversal alongside text and
embedding lookup.

\subsubsection{Graph-augmented retrieval}
\label{sec:lr-graph}

Graph-augmented methods supply the conceptual motivation for ingesting
OpenAlex into a typed knowledge graph but offer no benchmarked baseline
to copy on realistic scientific queries.  This thesis therefore treats
the literature as identifying the gap rather than filling it.  Scholarly
corpora are already graphs.  Papers cite papers, authors co-author,
venues publish, funders grant.  Three sub-threads have been proposed to
exploit this structure, namely knowledge-graph embeddings, graph neural
networks, and Personalized PageRank, but each has been evaluated almost
exclusively off scholarly search.

The canonical knowledge-graph embedding family treats link prediction as
a geometric problem.  TransE \cite{bordes2013transe}, ComplEx
\cite{trouillon2016complex}, and RotatE \cite{sun2019rotate} have been
evaluated almost exclusively on FB15k and WN18 and their harder variants,
benchmarks structurally unlike scholarly graphs in scale, edge
directionality, and node heterogeneity.  F\"arber and Ao
\cite{farber2022makgenhanced} apply vanilla KGEs at scholarly scale,
reporting ComplEx MRR 0.958 on paper, journal, and conference link
prediction over MAKG.  The literature does not, however, evaluate whether
such embeddings support natural-language literature search.
Link-prediction accuracy on the graph itself is not evidence of utility
for retrieval.

Graph neural networks face a similar mismatch.  GraphSAGE
\cite{hamilton2017graphsage} and GAT \cite{velickovic2018gat} are
evaluated on node classification over citation networks (Cora, Citeseer,
Pubmed) that are tiny next to OpenAlex, and the task is not retrieval.
LightGCN \cite{he2020lightgcn} simplified GCN-based collaborative
filtering on user-item bipartite graphs, but citation networks have no
recurring user, edges are directional and time-stamped, and the long
tail of once-cited papers is extreme.  The symmetric normalisation
assumptions LightGCN makes are not obviously valid in this setting.
Scholarly-specific neural work has largely bypassed these
bipartite-GNN assumptions \cite{bhagavatula2018citeomatic,shen2024tgntrec},
but none of it has produced a benchmark comparable in realism to
LitSearch.

Personalized PageRank has the strongest conceptual fit.  Haveliwala's
topic-sensitive formulation \cite{haveliwala2002tspr} encodes query
intent directly into the teleport vector, so the resulting scores reflect
relevance conditional on the query rather than the generic prominence of
raw citation counts, and the iterative propagation aggregates influence
across multiple citation hops.  Chen et al.\ \cite{chen2007gems} exploit
this to surface historically pivotal but under-cited papers in physics.
Together these properties address failure modes that plague both raw
citation ranking (no query conditioning) and dense retrieval (no
multi-hop aggregation over typed relations), making PPR an obvious
candidate for hybridisation with modern retrievers.  Despite this fit,
the 2020--2026 scholarly-search literature does not contain a rigorous
combination of Personalized PageRank with modern dense or KG embeddings
evaluated on scientific IR.  The closest analogue, HippoRAG
\cite{jimenezgutierrez2024hipporag}, builds a schemaless KG via LLM
OpenIE and seeds PPR from query entities, reporting Recall@5 of 51.9 on
MuSiQue and 89.1 on 2WikiMultiHopQA, but the corpora are Wikipedia-based,
not scientific.

Modern graph-RAG variants extend the framing further but evaluate
off-domain.  Microsoft's GraphRAG \cite{edge2024graphrag} clusters an
LLM-extracted graph with Leiden community detection and answers
sensemaking questions by aggregating community summaries, evaluated on
podcast transcripts and news articles.  LightRAG \cite{guo2024lightrag}
evaluates on synthetic sensemaking questions over textbook corpora.  For
scholarly search specifically, the strongest modern scholarly-RAG
system, OpenScholar \cite{asai2024openscholar}, abandons the citation
graph entirely and combines a text retriever with iterative
self-feedback.  The pattern is consistent.  Sophisticated graph methods
such as Personalized PageRank, GNNs, and community detection are widely
posited as useful for scholarly search but rarely benchmarked against
modern dense baselines on realistic scientific queries.

This thesis takes the conjunction of strong conceptual fit and thin
empirical record on scientific NL queries as motivation rather than as a
baseline to copy.  It ingests OpenAlex into a typed knowledge graph and
exposes graph structure through traversal queries directly, treating the
KGE, GNN and PPR literature as the work that identifies the gap rather
than the work that fills it.  The end-to-end generative component of
GraphRAG and its successors is out of scope for this thesis, since the
target is first-stage retrieval returning ranked paper lists rather than
synthesised long-form answers.  How users interact with that
graph-structured retrieval, in particular how natural-language questions
become structured queries against an OpenAlex-derived schema, is the
question Section~\ref{sec:lr-nl2q} takes up.

\subsubsection{Natural-language-to-structured-query systems for scholarly knowledge graphs}
\label{sec:lr-nl2q}

Structured-query translation is the design this thesis pursues,
instantiated through an intermediate JSON representation that decouples
LLM-based intent extraction from deterministic query compilation.  This
third position avoids the structural weaknesses of both existing
scholarly NL-to-graph-query benchmarks and end-to-end LLM tool calling.
The underlying technology has been maturing rapidly in the NL-to-SQL
setting.  Spider \cite{yu2018spider} and BIRD \cite{li2023bird}
established cross-domain benchmarks with strictly disjoint train/test
schemas, DIN-SQL \cite{pourreza2023dinsql} narrowed the human-model gap
through decomposition into schema linking, classification, generation
and self-correction, and Spider 2.0 \cite{lei2024spider2} drops GPT-4o to
roughly 10\% success on enterprise-scale schemas, showing how far
short-context benchmarks travel when realistic schemas appear.

The scholarly-specific literature is considerably thinner.  SpCQL
\cite{guo2022spcql} presents the first dedicated NL-to-Cypher dataset
with 10,000 question/query/answer triples in Chinese over the OwnThink
Neo4j graph, but the schema is not provided at inference, which leaves
the original setup under-specified.  DBLP-QuAD \cite{banerjee2023dblpquad}
builds 10,000 NL/SPARQL pairs over the DBLP scholarly graph using 98
hand-written templates expanded via subgraph sampling, and SciQA
\cite{auer2023sciqa} addresses SPARQL over the Open Research Knowledge
Graph with 2,565 pairs combining handcrafted and template-generated
queries.  Reported results on SciQA span from under 26\% F1 for zero-shot
LLMs to over 97\% F1 for fine-tuned 7--8B LLMs, a gap that is evidence
both for the importance of fine-tuning and for the risk of benchmark
saturation via template overfitting.  All three scholarly datasets share
structural weaknesses.  SpCQL is Chinese-only and open-domain, DBLP-QuAD
and SciQA are template-generated and reward systems that learn template
patterns rather than compositional semantics, and both SPARQL benchmarks
typically assume gold entity and relation linking that reality does not
supply.  QALD \cite{usbeck2023qald10} remains the general KGQA evaluation
of record but does not target scholarly graphs.

A very different line of work side-steps grammar-constrained semantic
parsing entirely.  ReAct \cite{yao2023react} interleaves free-text
thought traces with executable actions, and Toolformer
\cite{schick2023toolformer} shows an LM can self-supervise when to call
tools.  OpenAI's function-calling interface productised the pattern in
mid-2023 and is now the default substrate for agentic LLM systems.  The
contrast with fine-tuned parsers is structural.  Tool-using LLMs push
structured-query construction into the prompt and the tool's schema,
exchanging grammar-level correctness guarantees for flexibility and
compositional reasoning.  Direct head-to-head comparisons between
DIN-SQL-style fine-tuned parsers and ReAct-style tool-using agents on
scholarly KGQA are still scarce, leaving an open methodological question
particularly relevant to systems that expose OpenAlex through both a
Cypher backend and a set of graph-traversal functions.

This thesis's design occupies a third position.  Rather than fine-tuning
a grammar-constrained parser end-to-end (the SpCQL, DBLP-QuAD, and SciQA
tradition) or letting an LLM emit Cypher directly through function
calling (the ReAct and Toolformer tradition), it introduces an
intermediate JSON representation that names a target entity type and a
set of typed relations.  The LLM populates this structure from natural
language, and a deterministic compiler translates the JSON into Cypher
against a fixed schema derived from OpenAlex.  The split inherits
compositional flexibility from the function-calling line, since the LLM
does not need to know the underlying graph language, and preserves the
correctness guarantees of the fine-tuned parser line, since query
construction itself never leaves the compiler.  It also avoids three of
the empirical weaknesses surfaced above.  There are no templates to
memorise, gold entity linking is not assumed, and the schema is supplied
at inference rather than left implicit.  Whether this design generalises
to the kinds of structural queries reformulation cannot reach is the
question Section~\ref{sec:lr-rag} takes up.

\subsubsection{LLM query reformulation for scholarly search}
\label{sec:lr-rag}

This thesis sets LLM query reformulation aside, both because it degrades
the strong retrievers it would need to combine with and because it cannot
in principle express the structural queries a typed graph makes
available.  The relevant body of work is the cleanest basis for explaining
why.  A range of methods leave the retriever unchanged and rewrite the
query.  HyDE \cite{gao2023hyde} encodes an LLM-generated hypothetical
answer document, Query2doc \cite{wang2023query2doc} concatenates the
original query with an LLM pseudo-document, query-expansion-by-prompting
\cite{jagerman2023qexp} adds chain-of-thought expansion terms, and
Generative Relevance Feedback \cite{mackie2023grf} replaces classical
pseudo-relevance feedback with LLM-generated long-form text.

The critical corrective is Weller et al.\ \cite{weller2024whenfail}, who
ran eleven expansion methods across twelve datasets and twenty-four
retrievers and found a strong negative correlation between retriever base
strength and expansion gains.  Expansion helps weaker models like DPR and
untuned Contriever but routinely hurts ColBERTv2, SPLADEv2, GTE Large and
the E5 variants, with the proposed mechanism being that generated text
adds noise that dilutes the original relevance signal.  Two implications
follow.  A pipeline combining strong dense retrieval with a typed
OpenAlex-backed graph sits precisely where Weller et al.\ document
expansion as negligible or negative.  The structured graph is expected to
provide the discriminative signal expansion would otherwise contribute,
without the attendant noise.  Reformulation also cannot in principle
express structural queries such as \emph{papers by authors at institution
X that cite work from lab Y}, regardless of how semantically rich the
expansion.  The structured-parsing path surveyed in
Section~\ref{sec:lr-nl2q} is therefore the right tool for this thesis,
and reformulation is not.

A larger generative and agentic scholarly-RAG literature (PaperQA2
\cite{skarlinski2024paperqa2}, OpenScholar \cite{asai2024openscholar},
Ai2 Scholar QA, and commercial tools such as Elicit and Consensus) sits
adjacent to this thread but targets end-to-end answer generation rather
than first-stage retrieval.  Ranked paper lists are the deliverable here,
so natural-language question-answering is out of scope.

% Synthesising across the four threads, a coherent picture emerges for
% systems that combine OpenAlex-backed typed knowledge graphs with LLM-based
% query translation, text-index retrieval, graph traversal and dense
% embedding.  Dense retrievers are necessary but not sufficient.  The best
% open model on LitSearch still misses a quarter of relevant papers at
% Recall@5.  Citation-aware embeddings help paper--paper similarity more
% than natural-language ad-hoc search, and task-format adapters are almost
% certainly required.  Graph-augmented retrieval has strong conceptual
% motivation but a surprisingly thin evaluation record on scientific queries
% specifically, which is both a methodological risk and an opportunity.
% Natural-language-to-structured-query is feasible but its scholarly
% benchmarks are small, templated and approach saturation under fine-tuning.
% Hybrid approaches that combine LLM-based entity extraction with
% deterministic query compilation offer a less-studied alternative to both
% end-to-end semantic parsing and free-form LLM function-calling.  And LLM
% query reformulation is a powerful lever that improves weak pipelines but
% can degrade strong ones, making the relative position on the
% retriever-strength curve, rather than any absolute endorsement, the right
% frame for design decisions in hybrid scholarly search systems.